% =====================================================================
%  Chapter 19 — Online Bayesian Model Averaging
% =====================================================================

\chapter{Online Bayesian Model Averaging}\idx{Model averaging}
\label{ch:ensemble}

The book's model set has twelve members. This chapter explains how they
are combined, why the combination is done online rather than in one pass,
and why the combination is where the book's most important number is
produced.

\section{Why twelve, and not one}

Twelve members is three feature sets --- the full set, a trend-only set,
and a risk-only set --- crossed with two horizons and two prior
strengths. The crossing is not decoration. Each factor is a real
hypothesis:

\textbf{The feature set} is a hypothesis about where the signal is. The
full set says the signal is in the joint behaviour of price and volume.
The trend set says it is in the shape of the price path. The risk set says
it is in the risk characteristics of the fund. These are genuinely
different claims, and a book that tested only one of them would be making
a claim about the whole subject from one experiment.

\textbf{The horizon} is a hypothesis about the speed of the signal, as
Chapter 15 set out.

\textbf{The prior strength} is a hypothesis about how much the data should
be allowed to move the coefficients, and it interacts with the horizon: a
stronger prior matters more when the training set is small, which happens
at short horizons and in the early years.

Twelve members is also a *count*, and the count is a liability. Every
member is a trial, and Chapter 22's machinery exists to price that
liability. The book does not claim twelve is optimal; it claims that
twelve is the number of things actually tried, and that the number is
reported.

\section{Combining models online}

The combination is an online Bayesian model average. At each formation
date \(t\), every member produces a forecast, each member is scored
against the outcome that has just become known, and the members are
weighted in proportion to the exponential of their accumulated score.

The scoring rule for member \(m\) at date \(t\) uses the standardised
forecast error:
\begin{equation}
z_{i,m,t} = \frac{y_{i,t} - \mu_{i,m,t}}{\sqrt{v_{i,m,t}}},
\qquad
\ell_{m,t} = \frac{1}{N_t}\sum_{i} -\frac{1}{2}z_{i,m,t}^2 ,
\end{equation}
averaged over the cross-section. This is a proper scoring rule: a member
whose forecasts are calibrated scores about \(-1/2\) on average whether or
not it has skill, and a member that is systematically overconfident scores
worse than one that is appropriately uncertain. The rule therefore
rewards calibration first and skill second, which is the right order for
a system whose purpose is to size positions.

The horizon-aware part is the critical detail. A member with horizon \(h\)
produces a forecast at date \(s\) whose outcome is realised at \(s+h\). At
date \(t\), the score of that member uses the forecast made at \(t-h\),
because that is the forecast whose outcome is realised exactly now. The
book implements this by construction: the score is attached to the
formation date \(t-h\), and it becomes available to the ensemble\idx{Ensemble} only at
\(t\).

\begin{takeaway}[The rule that must not be broken]
A member's score at date \(t\) may use only an outcome that was realised
by \(t\). For a member with horizon \(h\), that means the forecast made at
\(t-h\). Using the forecast made at \(t\) and scoring it against the
outcome at \(t\) would let the ensemble weight itself on the future, and
the resulting equity curve would be a fiction.
\end{takeaway}

\section{Why the score is discounted}

The score is cumulative\idx{Cumulative}, and it is accumulated with an exponential discount:
\begin{equation}
L_{m,t} = \delta\, L_{m,t-1} + \ell_{m,t},
\qquad
\delta = 2^{-1/H},
\end{equation}
where \(H\) is a half-life\idx{Half-life} in months and the book uses \(H = 60\) by
default.

Two arguments justify the discount, and they are different.

The first is statistical. The score is an average over the cross-section
of a standardised error, and the cross-sectional average of a squared
standardised error has expectation one when the member is calibrated.
Summing it over hundreds of months accumulates the noise in that average
as if it were evidence, and the difference between two members' sums is
dominated by luck rather than by skill. The discount bounds the effective
sample size of the comparison.

The second is economic, and it is the more interesting one. Chapter 3
documented that published factor returns decay: the return to a signal
falls once it is widely known, because the arbitrage that pushed it up
exhausts itself. The same thing happens to a model built on price data.
A member that was excellent five years ago may be worthless now, and an
ensemble that remembers its whole record will weight it as though it were
still excellent. The half-life is a statement about how fast the book
believes a model's edge can change.

The half-life is a parameter, and parameters that can be chosen are
parameters that can be overfitted. It is therefore in the trial family:
the book runs the whole system with a half-life of twenty-four months and
with one of a hundred and twenty, alongside the default sixty, so that the
reader can see what the choice is worth.

\section{From the weights to the combined forecast}

With weights \(w_{m,t}\) summing to one, the combined posterior mean for
fund \(i\) is the weighted average of the members' means:
\begin{equation}
\mu_{i,t} = \sum_m w_{m,t}\,\mu_{i,m,t}.
\end{equation}

The combined posterior variance is where the interesting part is. By the
law of total variance,
\begin{equation}
v_{i,t} = \underbrace{\sum_m w_{m,t}\,v_{i,m,t}}_{\text{within}}
          + \underbrace{\sum_m w_{m,t}\left(\mu_{i,m,t} - \mu_{i,t}\right)^2}_{\text{between}} .
\end{equation}

The first term is the members' own uncertainty, averaged. The second is
the disagreement between the members, and it is the term that does the
architectural work of this book.

If the members agree, the between term is small, the combined forecast is
confident, and the portfolio layer is willing to tilt. If the members
disagree --- because one of them is reading the same price path as a trend
and another is reading it as risk --- the between term is large, the
combined forecast is uncertain, and the Kelly solve of Chapter 26
shrinks the position toward the stated policy.

The ensemble's own weights over the record are shown in \figref{fig:model-weights}. No configuration ever dominates, and the weights move slowly.

This is the sentence the book has been building towards: \emph{the
machine-learning layer's real job is to estimate the posterior over the
edge, and the disagreement between models is itself information about how
much edge there is.} A model that reported only a point forecast and no
uncertainty would be a different system, and a worse one --- not because
it would be wrong, but because it would not know that it was wrong.

\begin{keytakeaways}
\item Twelve members, three feature sets by two horizons by two prior
strengths, each a real hypothesis and each a trial whose count is
published.
\item The scoring rule is a proper scoring rule on the standardised
forecast error, applied horizon-aware so that no member is ever scored on
an outcome that had not yet been realised.
\item The score is discounted with a sixty-month half-life, because
accumulated noise otherwise masquerades as skill and because the edge
itself decays.
\item The combined variance carries a between-model term, and that term is
what makes the portfolio shrink toward the policy when the members
disagree.
\end{keytakeaways}
